% ═══════════════════════════════════════════════════════════════
% ML-07a: Musterlösung - Comida (Lebensmittel)
% Abgeleitet aus LE-07a (Abschnitt 9: Lösungen)
% ═══════════════════════════════════════════════════════════════

\documentclass[11pt, a4paper]{article}

\usepackage[utf8]{inputenc}
\usepackage[T1]{fontenc}
\usepackage[ngerman]{babel}
\usepackage{helvet}
\renewcommand{\familydefault}{\sfdefault}

\usepackage[margin=2cm]{geometry}
\usepackage{enumitem}
\usepackage{tabularx}
\usepackage{booktabs}
\usepackage{multicol}

\usepackage{xcolor}
\usepackage{tcolorbox}

\usepackage{fancyhdr}
\usepackage{lastpage}

% Farben
\definecolor{spanisch}{HTML}{c41e3a}
\definecolor{grau}{HTML}{888888}
\definecolor{hellgrau}{HTML}{f5f5f5}
\definecolor{orange}{HTML}{b35c00}
\definecolor{loesung}{HTML}{c62828}

\newcommand{\fachfarbe}{spanisch}

% Variablen
\newcommand{\thema}{Comida -- Lebensmittel}
\newcommand{\sequenz}{07}
\newcommand{\block}{a}
\newcommand{\fach}{Spanisch}
\newcommand{\klasse}{7}
\newcommand{\maxpunkte}{24}

% Kopf-/Fußzeile
\pagestyle{fancy}
\fancyhf{}
\renewcommand{\headrulewidth}{0.4pt}
\renewcommand{\footrulewidth}{0.4pt}
\lhead{\fach{} -- Klasse \klasse}
\chead{\textcolor{loesung}{\textbf{ML-\sequenz\block{} (MUSTERLÖSUNG)}}}
\rhead{}
\lfoot{}
\rfoot{Seite \thepage\ von \pageref{LastPage}}

% Befehle
\newcommand{\punktebox}[1]{\hfill\fbox{\small \_/#1}}
\newcommand{\loesung}[1]{\textcolor{loesung}{\textbf{#1}}}
\newcommand{\luecke}[1]{\rule{#1}{0.4pt}}
\newcommand{\es}[1]{\textcolor{\fachfarbe}{\textbf{#1}}}

% Merkkasten
\newtcolorbox{merkkasten}[1][]{
    colback=hellgrau,
    colframe=\fachfarbe,
    boxrule=1pt,
    arc=0pt,
    left=8pt, right=8pt, top=8pt, bottom=8pt,
    fonttitle=\bfseries\color{\fachfarbe},
    title=#1
}

% Hinweis
\newtcolorbox{hinweis}{
    colback=white,
    colframe=orange,
    boxrule=1pt,
    arc=0pt,
    left=5pt, right=5pt, top=5pt, bottom=5pt,
    fonttitle=\bfseries,
    title=Hinweis
}

% Aufgabe
\newenvironment{aufgabe}[1]{%
    \vspace{10pt}
    \noindent\textbf{\textcolor{\fachfarbe}{Aufgabe #1}}
    \vspace{5pt}
    \begin{list}{}{\leftmargin=0pt}
    \item[]
}{%
    \end{list}
}

% ═══════════════════════════════════════════════════════════════
\begin{document}

% Kopfbereich
\begin{center}
    {\Large\bfseries\textcolor{\fachfarbe}{Arbeitsblatt: \thema}}\\[0.3em]
    {\large\textcolor{loesung}{\textbf{--- MUSTERLÖSUNG ---}}}
\end{center}

\vspace{5pt}

\noindent
\begin{tabularx}{\textwidth}{|X|X|X|}
\hline
\textbf{Name:} \textcolor{loesung}{Musterschüler/in} &
\textbf{Klasse:} \klasse &
\textbf{Datum:} \textcolor{loesung}{XX.XX.XXXX} \\
\hline
\end{tabularx}

\vspace{15pt}

% ═══════════════════════════════════════════════════════════════
% MERKKASTEN (mit Lösungen)
% ═══════════════════════════════════════════════════════════════

\begin{merkkasten}[Merkkasten: Lebensmittel (La comida)]

\begin{tabularx}{\textwidth}{|l|X|}
\hline
\textbf{La fruta} (Obst) & \loesung{la manzana, el plátano, la naranja} \\[0.5em]
\hline
\textbf{La verdura} (Gemüse) & \loesung{el tomate, la lechuga, la zanahoria, la cebolla} \\[0.5em]
\hline
\textbf{La carne / El pescado} & \loesung{la carne, el pollo, el pescado} \\[0.5em]
\hline
\textbf{Otros} (Andere) & \loesung{el pan, la leche, el queso, el huevo, el arroz} \\[0.5em]
\hline
\end{tabularx}

\end{merkkasten}

\vspace{10pt}

% ═══════════════════════════════════════════════════════════════
% AUFGABE 1 (mit Lösungen)
% ═══════════════════════════════════════════════════════════════

\begin{aufgabe}{1}
Verbinde die deutschen Wörter mit dem spanischen Begriff. \punktebox{5}
\end{aufgabe}

\begin{center}
\begin{tabular}{rcl}
Apfel & \loesung{------} & \loesung{la manzana} \\[0.8em]
Banane & \loesung{------} & \loesung{el plátano} \\[0.8em]
Tomate & \loesung{------} & \loesung{el tomate} \\[0.8em]
Salat & \loesung{------} & \loesung{la lechuga} \\[0.8em]
Milch & \loesung{------} & \loesung{la leche} \\
\end{tabular}
\end{center}

\textit{\textcolor{grau}{Je 1P pro richtige Zuordnung}}

\vspace{15pt}

% ═══════════════════════════════════════════════════════════════
% AUFGABE 2 (mit Lösungen)
% ═══════════════════════════════════════════════════════════════

\begin{aufgabe}{2}
Ordne die Lebensmittel in die richtige Kategorie ein. \punktebox{6}
\end{aufgabe}

\begin{tabularx}{\textwidth}{|X|X|}
\hline
\textbf{La fruta} & \textbf{La verdura} \\
\loesung{la naranja, el plátano} & \loesung{la zanahoria, la cebolla} \\[1em]
\hline
\textbf{La carne / El pescado} & \textbf{Otros} \\
\loesung{el pollo, el pescado, la carne} & \loesung{el queso, el pan, la leche} \\[1em]
\hline
\end{tabularx}

\vspace{0.5em}
\textit{\textcolor{grau}{Je 0.5P pro richtiger Zuordnung, max. 6P}}

\vspace{15pt}

% ═══════════════════════════════════════════════════════════════
% AUFGABE 3 (mit Lösungen)
% ═══════════════════════════════════════════════════════════════

\begin{aufgabe}{3}
Setze die richtige Mengenangabe ein. \punktebox{5}
\end{aufgabe}

\begin{enumerate}[label=\alph*)]
    \item Quiero \loesung{un kilo de} manzanas.
    \item Necesito \loesung{medio kilo de} tomates.
    \item Dame \loesung{un litro de} leche.
    \item Quiero \loesung{una botella de} agua.
    \item Necesito \loesung{un paquete de} arroz.
\end{enumerate}

\textit{\textcolor{grau}{Je 1P pro richtige Mengenangabe}}

\vspace{15pt}

% ═══════════════════════════════════════════════════════════════
% MERKKASTEN: Mengenangaben (ausgefüllt)
% ═══════════════════════════════════════════════════════════════

\begin{merkkasten}[Merkkasten: Mengenangaben (Cantidades)]

\begin{tabular}{ll}
\es{un kilo de...} & ein Kilo... \\
\es{medio kilo de...} & ein halbes Kilo... \\
\es{un litro de...} & ein Liter... \\
\es{una botella de...} & eine Flasche... \\
\es{un paquete de...} & ein Paket/eine Packung... \\
\end{tabular}

\vspace{0.5em}
\textbf{Wichtig:} Nach Mengenangaben steht immer \es{de}!
\end{merkkasten}

\vspace{15pt}

% ═══════════════════════════════════════════════════════════════
% AUFGABE 4 (Beispiellösung)
% ═══════════════════════════════════════════════════════════════

\begin{aufgabe}{4}
Schreibe eine Einkaufsliste auf Spanisch (mindestens 5 Produkte mit Mengenangaben). \punktebox{8}
\end{aufgabe}

\textbf{Beispiellösung -- Mi lista de compras:}

\begin{enumerate}
    \item \loesung{un kilo de manzanas}
    \item \loesung{medio kilo de tomates}
    \item \loesung{un litro de leche}
    \item \loesung{un paquete de arroz}
    \item \loesung{una botella de agua}
    \item \loesung{200 gramos de queso} \textit{(Bonus)}
\end{enumerate}

\vspace{0.5em}
\textit{\textcolor{grau}{Bewertung:}}
\begin{itemize}
    \item Mengenangabe korrekt: je 1P (max. 5P)
    \item Lebensmittel korrekt (Artikel + Wort): je 0.5P (max. 2.5P)
    \item Bonus (6. Produkt): +0.5P
\end{itemize}

\vspace{20pt}
\hrule
\vspace{10pt}

\begin{flushright}
\textbf{Punkte:} \loesung{\maxpunkte} / \maxpunkte
\end{flushright}

\end{document}
